\section{AI for Math 技术栈：从直觉到可验证证明}

很多人把 AI for Math 简化成“模型会不会做数学题”。这只是最浅的一层。真正的目标，是让 AI 从问题、猜想、非形式推理出发，逐步走向形式化陈述、证明搜索、机器验证和人类反馈。

\begin{figure}[H]
\centering
\includegraphics[width=0.94\textwidth]{ai-for-math-stack.png}
\caption{AI for Math 技术栈：从问题到验证再到数学家反馈。}
\end{figure}

\begin{knowledgebox}{读图：AI for Math 不是单点模型能力}
图中每一层都可能成为瓶颈：问题理解需要语义和背景知识；非形式推理需要直觉；形式化需要 Lean/Coq 语言；证明搜索需要定理库和 tactic；验证需要 proof assistant；最后还需要数学家判断方向和意义。
\end{knowledgebox}

\subsection{术语消化：AI for Math 关键词}

\noindent\begin{tabularx}{\textwidth}{p{0.22\textwidth}p{0.32\textwidth}X}
\toprule
术语 & 解决的问题 & 与本期关系 \\
\midrule
Lean & 形式化数学语言和 proof assistant & 把自然语言证明转成机器可检查证明。 \\
Coq & 另一类 proof assistant & 说明形式化证明不止 Lean 一条路线。 \\
Tactic & 自动或半自动证明步骤 & 让模型能逐步构造可检查证明。 \\
Theorem Proving & 定理证明 & AI for Math 的核心任务之一。 \\
Verifier & 验证器 & 判断证明是否合法，避免“看起来对”的幻觉。 \\
Formalization & 形式化 & 把论文、定义、猜想写进可验证系统。 \\
\bottomrule
\end{tabularx}

\subsection{为什么 proof assistant 是关键}

数学自然语言证明经常省略步骤，依赖读者背景和直觉。Proof assistant 要求每一步都符合规则。对 AI 来说，这是一种强反馈：生成的证明要么被 kernel 接受，要么失败并给出目标状态。这种反馈比主观评分更硬，也更适合训练和评估。

\begin{importantbox}{可验证性是 AI for Math 的护城河}
数学是少数能提供强验证信号的高价值领域。Lean/Coq 等系统让模型输出不只是“像证明”，而是必须通过机器检查。这使 AI for Math 同时具备科学价值和工程可训练性。
\end{importantbox}

\subsection{本章小结}

AI for Math 的技术栈包括直觉、形式语言、证明搜索和机器验证。真正困难的是把这些层连起来，而不是只在某一层做题。

